\section{The product: a shared specification that can be reviewed and executed}
\label{sec:product}

LegalMath should be a regulatory specification workbench. Its primary users are
the compliance officer interpreting a change, the product or operations owner
deciding how it affects the business, and the engineer implementing the resulting
control. An auditor or independent reviewer should be able to reconstruct the
same decision later. The product succeeds when those people can inspect one
maintained interpretation through views suited to their work.

The first release should support a narrow end-to-end task: take a selected public
circular and its dependencies, prepare proposed rules, resolve or retain
interpretation questions, evaluate synthetic cases, and export a versioned change
package. Its result is a reviewed specification with evidence. Integration into
bank systems follows the bank's ordinary change process.

\subsection{What a reviewer should see}

Opening the SPI financial condition should display the source passage and its
annex context beside a controlled explanation: the portfolio route or the
net-asset route can satisfy this particular condition. Definitions should be
available in place, especially ownership attribution and exclusion of the primary
residence. The same screen should show the boundary cases, missing evidence and
where the condition is used in the larger SPI assessment.

The reviewer can change an interpretation, ask for a distinguishing example, or
record a question for the policy owner. If two candidate readings differ only at
the threshold, the workbench should produce that boundary case. If they differ
because one applies to a representative and another to the client, the example
should name the roles. The user should not need to understand a solver formula
to see why the choice matters.

The engineer views the corresponding types, data mappings, executable dependency
graph and tests. The operations owner views required actions, responsible teams,
timing and evidence of completion. A diagram or decision table is generated from
the same adopted rules. Maintaining a separate hand-edited flowchart would create
another translation boundary and another version to reconcile.

\begin{figure}[htbp]
\centering
\begin{tikzpicture}[node distance=8mm and 8mm,
 box/.style={draw=teal,fill=pale,rounded corners,text width=40mm,align=center,minimum height=17mm,font=\small},
 arr/.style={-{Latex[length=2mm]},thick,draw=navy}]
\node[box] (source) {Preserved sources\\and dependencies};
\node[box,right=of source] (draft) {Proposed interpretation\\and unresolved questions};
\node[box,right=of draft] (spec) {Reviewed typed rules\\and examples};
\node[box,below=of spec] (exec) {Interpreter and\\integration adapters};
\node[box,left=of exec] (evidence) {Tests, counterexamples\\and scoped proofs};
\node[box,left=of evidence] (history) {Decision explanations\\and amendment history};
\draw[arr] (source)--(draft); \draw[arr] (draft)--(spec);
\draw[arr] (spec)--(exec); \draw[arr] (exec)--(evidence);
\draw[arr] (evidence)--(history);
\draw[arr] (evidence.north)--(draft.south);
\end{tikzpicture}
\caption{The proposed workbench. Counterexamples return to interpretation review.
A proof result remains attached to the exact rules, assumptions and software
version it checks.}
\label{fig:architecture}
\end{figure}

\subsection{The intermediate representation}

A JSON document becomes an executable specification only when its expressions
have defined meanings. RuleIR 0.1 fixes a finite, acyclic expression language and
a separate event profile. The complete grammar is supplied as JSON Schema;
\path{docs/specs/v0.1/semantics.md} specifies evaluation; and
\path{docs/specs/v0.1/contracts.md} specifies identity, storage and service
behavior. An implementer can start from these contracts without selecting a
logic engine or inventing missing-value behavior.

Each rule needs more than its condition. It must identify the source document and
revision, paragraph or page, exact supporting span, language, source digest and
retrieval date. It must record the legal actor, activity, product, client class,
relevant time and interpretation status. It must distinguish a derived fact from
a duty, permission or prohibition. A duty also needs an actor, trigger, action,
timing rule and evidence by which performance can be assessed. A missing source
deadline remains missing; operational escalation times belong to a separately
identified bank policy.

Definitions should be reusable and versioned. A field for portfolio value must
refer to the adopted portfolio definition, relevant ownership rules, valuation
date and currency conversion convention. The application mapping records where
each fact comes from, its freshness and who can resolve a conflict. A rule whose
necessary definition remains unresolved can be reviewed and discussed, but it
cannot be released as an executable control for that use.

The complete financial-condition bundle is supplied in
\path{fixtures/spi-financial.bundle.json} beneath the specification directory.
It declares two money inputs, an unconditional scope for this isolated
subcondition, the expression, interpretation record and a verified span in
Annex 1 paragraph 3.1. The excerpt below is a complete expression node from
that bundle. Node identifiers make its trace reproducible.

\begin{lstlisting}[caption={The portfolio comparison in the supplied JSON bundle. Amounts are HK cents.}]
{"node_id":"portfolio.ge", "op":"compare", "cmp":"ge",
 "left":{"node_id":"portfolio.input", "op":"fact",
         "name":"portfolio"},
 "right":{"node_id":"portfolio.threshold", "op":"literal",
          "type":"money_hkd", "value":"4000000000"}}
\end{lstlisting}

The second comparison uses \texttt{net\_assets\_ex\_home} and
\texttt{"8000000000"}; an \texttt{any} node combines them. These inputs denote
amounts normalized under adopted wealth definitions, not any account balance
supplied under the same name. Release requires those definitions and applicability
to be reviewed. The narrow result name, \texttt{spi.financial}, identifies a
subcondition rather than a complete SPI assessment or trade authorization.

\subsubsection{Types and the expressions an agent may generate}

The scalar types are Boolean, integer, HKD money and Gregorian date. Integer
and money values use canonical decimal strings; monetary units are cents.
There are no floating-point values, implicit casts or known nulls. A date must
be valid, including its leap-year constraints. Foreign-currency normalization
records the rate, valuation instant and rounding policy before producing an HKD
input. A missing exchange rate therefore leaves an unknown amount.

The expression grammar contains typed literals, fact and rule references,
\texttt{all}, \texttt{any}, \texttt{not}, equality/inclusive/strict comparison,
addition, subtraction, rational scaling, conditional choice and explicit defaults.
Addition and subtraction require matching integer or money types. Scaling uses
a nonnegative integer numerator and a positive denominator; a known result must
divide exactly. A half-cent allocation returns \texttt{E\_INEXACT\_SCALE} until
an external reviewed rounding policy resolves it. Date arithmetic is kept in
the separate scheduler so ordinary library behavior cannot choose a legal
anniversary convention unnoticed.

Every node has a unique identifier. References must resolve, types must agree
and the complete dependency graph must be acyclic. A rule reference may refer
only to an unconditional definition; a scoped result cannot quietly become a
Boolean false in another expression. Authors can expose the scope explicitly
as a named Boolean definition. Unknown operators and additional JSON properties
are rejected. This keeps a model from extending the language merely by emitting
a plausible field name.

\subsubsection{Evidence is an input to computation, with its own history}

A normalized fact is known, unknown or conflicting. Known facts have an exact
value, evidence identifiers, a validity interval and a recorded timestamp.
Unknown facts explain whether evidence is missing, stale or awaiting assessment.
Conflicts preserve the disagreeing records. Two admissible records with the
same value can combine their evidence; different values require resolution.
Arrival order is not a ranking of legal authority.

Every request supplies \texttt{valid\_at}, the instant being assessed, and
\texttt{known\_at}, the latest evidence the assessment may use. Intervals are
half-open: an amount valid until midnight is unavailable at that midnight.
A correction names the record it supersedes and becomes visible only when its
recorded time is eligible. Thus a corrected assessment can use new information
while the original decision remains reproducible.

Before Boolean evaluation, the engine checks conflicts in the static transitive
dependency closure, including branches that might later be skipped. A conflicting
portfolio amount therefore returns \texttt{CONFLICT} even if the net-asset route
could establish the financial condition. An unrelated conflict does not block
that rule. This conservative policy is a proposed engineering choice; it is
deliberately different from ordinary unknown-value propagation.

\subsubsection{Defaults, branch selection and failure behavior}

A default contains a base value and named exceptions. Each exception has a
Boolean guard, a value, an interpretation identifier and source spans. Evaluate
all guards first. Two true guards produce a conflict, even if the proposed values
agree. Otherwise any unknown guard yields an unknown result, because it might
become a competing exception. With exactly one true guard and all others false,
evaluate that exception's value. With all false, evaluate the base. Priority is
expressed only through explicitly reviewed nesting, never row order.

Boolean combinations evaluate all operands in source order. A runtime error or
exception conflict cannot disappear because a different operand is sufficient.
A conditional evaluates its condition and then only the selected branch; an
unknown condition yields unknown, even when both branch values happen to be
equal. Skipped branches are still type checked. These choices eliminate
implementation-dependent short-circuit behavior while retaining precise traces.

The rule's scope is evaluated separately. A false scope yields
\texttt{OUT\_OF\_SCOPE}; an unknown scope yields \texttt{UNKNOWN}. In scope,
Boolean bodies return \texttt{TRUE} or \texttt{FALSE}, other known bodies return
\texttt{VALUE}, and missing evidence returns \texttt{UNKNOWN}. Conflicts and
errors are distinct tagged results. An invalid version time is an error, not a
negative assessment. No exception handler may fall back to false.

The response carries all encountered missing inputs and the conservative subset
blocking the answer. If a known net-asset route establishes the alternative,
the unknown portfolio remains visible but is not a blocker. For an unknown
result, the union of encountered unknown dependencies is reported; the prototype
does not claim it is a mathematically minimal explanation. Traces record evaluated
nodes in postorder, evidence, source spans and explicitly skipped branch roots.
Explanations and diagrams are rendered from that trace.


\subsection{A precise example: money, alternatives and unknown values}
\label{sec:semantics}

Let $P$ denote the client's portfolio amount under the adopted definition, and
let $N$ denote net assets excluding the primary residence. Both are valued at the
relevant date and expressed in HK dollars under an explicit conversion policy.
For complete, valid inputs, the financial condition is
\begin{equation}
 F(P,N)=(P\geq 40{,}000{,}000)\ \lor\ (N\geq 80{,}000{,}000).
 \label{eq:financial}
\end{equation}
The source of the constants, inclusive comparisons and alternative is Annex 1,
paragraph 3.1. Equation~\eqref{eq:financial} formalizes that subcondition; its
inputs incorporate the separate definitions in paragraphs 3.2--3.4
\citep{sfcspiannex1}. RuleIR stores monetary amounts as
integer HK cents encoded as decimal strings. Conversions must take place
before the comparison, with the exchange-rate source and valuation instant
recorded. Floating-point rounding must not silently move a boundary.

Bank evidence is often incomplete. Define the comparison with a missing amount
to return $\Unknown$, and let $\True$, $\False$ and $\Unknown$ mean established,
refuted and unresolved for this condition. For the alternative in
Equation~\eqref{eq:financial}, use the following truth table:
\begin{equation}
\begin{array}{c|ccc}
\lor & \True & \False & \Unknown\\\hline
\True & \True & \True & \True\\
\False & \True & \False & \Unknown\\
\Unknown & \True & \Unknown & \Unknown
\end{array}
\label{eq:or}
\end{equation}
This table follows the intended treatment of incomplete evidence: a known
sufficient route establishes the alternative, both refuted routes refute it,
and otherwise missing information can matter. With unknown $P$ and
$N=90{,}000{,}000$, the second comparison is true, so
$\Unknown\lor\True=\True$. With unknown $P$ and $N=70{,}000{,}000$, the result
is $\Unknown\lor\False=\Unknown$. Figure~\ref{fig:financial-flow} gives a
flowchart view.

\begin{figure}[!htbp]
\centering
\begin{tikzpicture}[font=\small,
 box/.style={draw=teal,fill=pale,rounded corners,text width=39mm,align=center,minimum height=13mm},
 decision/.style={draw=teal,diamond,aspect=2.6,align=center,text width=40mm,inner sep=1mm},
 arr/.style={-{Latex[length=2mm]},thick,draw=navy}]
\node[box] (input) at (0,0) {Defined, dated amounts\\or explicit unknowns};
\node[decision] (yes) at (0,-2.8) {Is either financial\\route established?};
\node[box] (true) at (6,-2.8) {Financial condition\\established};
\node[decision] (no) at (0,-5.6) {Are both financial\\routes refuted?};
\node[box] (false) at (6,-5.6) {Financial condition\\refuted};
\node[box] (unknown) at (0,-8.3) {Unresolved: identify\\decision-relevant evidence};
\draw[arr] (input)--(yes);
\draw[arr] (yes)--node[above]{Yes}(true);
\draw[arr] (yes)--node[left]{No}(no);
\draw[arr] (no)--node[above]{Yes}(false);
\draw[arr] (no)--node[left]{No}(unknown);
\end{tikzpicture}
\caption{Proposed review flow for the financial subcondition, after resolving
input conflicts. Passing this test alone does not establish SPI status.}
\label{fig:financial-flow}
\end{figure}

The other Boolean operations need equally explicit tables. Under the proposed
conservative three-valued semantics, conjunction is false if an operand is false,
true if both are true, and otherwise unresolved; negation swaps true and false
and preserves unknown. This semantics does not discover every conclusion implied
by correlations among unknown facts. For example, a general expression
$A\lor\neg A$ remains unresolved when $A$ is unresolved under these tables.
A more powerful completion-based analysis may establish additional conclusions,
but it must be a specified refinement with its own evidence. No accidental
change of logic should occur when switching backends.

Contradictory input records are handled before this evaluation. If two
authoritative-looking records disagree on the same ownership allocation, the
workbench creates an evidence conflict and identifies the disputed fact. It must
not insert both as classical premises and infer arbitrary results from an
inconsistent theory. Being outside a rule's scope, having insufficient evidence,
and detecting a data conflict are also distinct operational outcomes. These
statuses should wrap the logical result rather than being encoded as false.

\subsection{Obligations and event histories}

The decision fragment cannot represent every continuing control. For consent,
let $s_i$ describe the recorded arrangement immediately before event $e_i$:
its client, selected categories, agreed threshold, monitoring mode, relevant
reviews and withdrawal status. Each event includes its timestamps and ordering
evidence. It changes the recorded state through a
versioned transition function,
\begin{equation}
 s_{i+1}=T_v(s_i,e_i).
 \label{eq:transition}
\end{equation}
Here $i$ orders processed events; it does not assume equal physical time
steps. Version $v$ identifies the adopted rules and ordering convention. A
withdrawal event can change the availability of the streamlined arrangement for
subsequent decisions. A new-funds event can trigger the designated-account
checks. An annual-review event creates work that must later be evidenced as
completed. The exact consequences require review of the cited guidance and the
bank's operating procedure \citep{sfcspiannex1,sfcspiannex2}.

An obligation record should distinguish its creation, due condition, performance,
breach and any later remediation. If a transaction occurs after a withdrawal, the
audit view must retain the transaction and identify the issue; it cannot erase
the event because the authorization view would have blocked it. This follows the
useful separation of normative permission from actual occurrence in eFLINT
\citep{eflint2026}.

Time requires two histories. Legal validity time records when a source or adopted
rule applies. System-recorded time records when the bank received and deployed
it. These support different questions: what should have applied to an event, and
what version did the system actually use? The same distinction applies to facts
received late or corrected later. A reconstructed historical assessment may use
new evidence, but it must not overwrite the original decision and its inputs.

Event ordering is a material assumption. The prototype should store source
timestamps, receipt timestamps, sequence identifiers, time zone and any uncertainty
in ordering. A synthetic trace with a consent withdrawal and execution at the
same timestamp is unresolved unless the available ordering or an adopted
convention settles it. Calendar computations should identify whether they use
anniversaries, elapsed days or another reviewed operation, drawing on the date
semantics literature \citep{dates2024}. Neither the operating system clock nor the
order in which messages happen to arrive establishes legal priority.

\subsubsection{The finite event profile}

The first event implementation supports consent and non-preemptive achievement
duties. Events have immutable identifiers, stream identifiers, occurrence and
recorded timestamps, and an authoritative sequence number. Timestamps are UTC
with microsecond precision; legal date conversion explicitly uses
\texttt{Asia/Hong\_Kong}. Equal occurrence times can use sequence order only
when the connector declares that ordering authority. Otherwise the result is
\texttt{E\_ORDER\_UNRESOLVED}. A duplicate event with identical bytes is a
no-op; reuse of its identifier for different content is an error.

Consent begins as not established only when the stream is complete from its
declared inception, or is complete after a reviewed initial snapshot. An incomplete export
instead yields an unknown history. A grant creates an active consent generation;
withdrawal makes it withdrawn; a new grant needs new evidence and creates a new
generation. Observed transactions remain in the audit history even when the
current consent state would prevent streamlined use.

An obligation records actor, action, subject, source version, activation and a
deadline or explicit absence of a deadline. Let $a$ be activation and $d$ the
exclusive deadline. Under this proposed profile, a matching performance at time
$p$ is timely when
\begin{equation}
 a\leq p<d.
 \label{eq:timely}
\end{equation}
This is the chosen boundary convention for the profile, not an assertion that
every SFC deadline has the same wording. A local due date is converted to the
next local midnight. An anniversary requires an explicit leap-day policy; a
reviewed choice of February 28, March 1 or rejection replaces an implicit library
default. Business-day calendars are deferred.

Absence of a timely record establishes breach only when a watermark shows that
the observation window through $d$ is complete. The wall clock cannot create
that evidence. Without completeness the obligation stays open. Timely performance
produces \texttt{SATISFIED\_ON\_TIME}; a complete window without it produces
\texttt{BREACHED}. A subsequent matching performance produces
\texttt{PERFORMED\_LATE} and retains the breach. A deadline-free request can be
performed but does not become overdue through an invented number of days.

For example, consider a synthetic review activated at midnight UTC on September
1 and due at midnight on September 2. A record one microsecond before the latter
instant is timely; a record exactly at it is late. The fixtures contain both.
These dates illustrate the runtime convention rather than an SFC-prescribed
review deadline. If a late-arriving record later proves timely performance,
replay produces a corrected assessment with a new knowledge cutoff; the original
assessment remains visible. Maintenance duties, pre-emptive satisfaction,
waivers and legal compensation require future profiles and are rejected now.


\subsection{What can be proved, and what must be reviewed}
\label{sec:proof}

There are three distinct assurance questions. The first is interpretive: does the
adopted specification represent the relevant sources for the chosen business
scope? The second is computational: does the implementation compute that
specification? The third is evidential: do the facts and event records represent
the relevant circumstances? A complete change package needs evidence for all
three, but the evidence is different.

To make the computational question precise, let $R_v$ be an approved version of
the formal rules. Let $\llbracket R_v\rrbracket(x)$ denote the reference meaning
of those rules on input $x$, and let $C_v(x)$ be the result of an implementation
or compiled adapter. Define $D_v$ to include the declared input types, admissible
values, explicit unknowns and relevant state. The target is
\begin{equation}
 \forall x\in D_v:\quad C_v(x)=\llbracket R_v\rrbracket(x).
 \label{eq:refinement}
\end{equation}
Equality here includes the named decision and required obligations represented by
the specification. If explanations are part of the verified output, their
relationship must be stated too. The meaning of the output cannot shrink from a
set of duties to one Boolean while retaining the same assurance claim.

Equation~\eqref{eq:refinement} is a proposed verification objective, not a result
already established by LegalMath. For a finite input space it may be checked by
enumeration; for a restricted expression language it may be amenable to an SMT
encoding or a proof by induction over expressions. For a production adapter,
the mapping from bank data into $x$ and the implementation of calendar and
currency operations also belong to the argument. A test suite samples the
equation and provides useful defect evidence, but it does not generally prove its
universal quantifier.

The interpretive question does not become a theorem simply by naming a formal
target. Source text, legal context and justified interpretations determine
$R_v$. Compliance review, source coverage, competing readings and independently
constructed cases address that step. Likewise, a theorem cannot establish that
a recorded balance or signature is authentic. These limits are operationally
useful: each review task is directed to the person and evidence that can answer it.

\subsection{Counterexamples make disagreements reviewable}

For many regulatory changes, a small counterexample is more useful than a large
proof transcript. Suppose the proposed implementation accidentally replaces the
first inclusive comparison in Equation~\eqref{eq:financial} with a strict one.
Let that candidate be $G$, and write $D$ for the declared set of admissible,
complete inputs $(P,N)$. A solver can search for
\begin{equation}
 \exists (P,N)\in D:\quad F(P,N)\ne G(P,N).
 \label{eq:counterexample}
\end{equation}
The concrete values $P=40{,}000{,}000$ and $N=0$ witness the disagreement:
$F$ is true and $G$ is false. A compliance officer can understand the error
without reading the encoding. Replacing the alternative with a conjunction is
similarly exposed by $P=35{,}000{,}000$ and $N=90{,}000{,}000$.

This technique also compares two candidate interpretations or old and new rule
versions. It must check the admissibility and consistency of its assumptions
first. If a test domain accidentally demands that the same amount is both
positive and negative, the absence of a counterexample proves nothing useful
about real cases. Unsatisfiable assumptions should be reported as a separate
failure, with the conflicting constraints identified. Solver timeout and unknown
results should never be rendered as ``equivalent.''

The initial Z3 adapter supports a smaller language than the interpreter:
Boolean, integer and money expressions with explicit knownness and bounded
declared domains. Dates, scaling, defaults and events return
\texttt{UNSUPPORTED}. Encode a Boolean expression $a$ by two bits $(t_a,f_a)$:
true is $(1,0)$, false $(0,1)$ and unknown $(0,0)$; prohibit $(1,1)$.
Conjunction then has the encoding
\begin{equation}
 t_{a\land b}=t_a\land t_b,\qquad
 f_{a\land b}=f_a\lor f_b.
 \label{eq:knownness}
\end{equation}
Disjunction uses the dual equations and negation swaps the bits. A numeric value
carries a knownness bit, and a comparison produces neither Boolean bit unless
both operands are known. These equations implement the stated truth tables.
Giving an unknown fact an unconstrained classical Boolean would instead make
$A\lor\neg A$ always true and verify a different language.

The service checks domain satisfiability before searching for a difference in
tagged results. Display a satisfying model only after replay through both
interpreters. An unsatisfiable difference query yields
\texttt{NO\_COUNTEREXAMPLE\_IN\_DECLARED\_DOMAIN}; an empty domain yields
\texttt{INCONSISTENT\_DOMAIN}; timeout stays \texttt{UNKNOWN}. Preserve formula,
domain, rule hashes and solver version. A solver answer is a different class
of evidence from a certificate replayed by an independent checker.

For event rules, the corresponding search concerns traces. A useful first
property is that, in the adopted consent model, no transition representing a new
streamlined decision uses a withdrawn consent state. Another is that processing
the same evidence event twice does not create duplicate obligations. Each
property has assumptions about event identity and ordering. Bounded trace search
can find violations, while a proof of an invariant can establish a stronger
claim if the transition model and induction obligations are checked.

\subsection{The smallest useful execution architecture}

The first runtime is a Python 3.11 reference interpreter, with JSON Schema and
semantic validation, SQLite storage, content-addressed source files and a local
FastAPI service. A server-rendered interface presents source, rule and case
together. The interpreter walks a validated syntax tree; it does not execute
model-generated Python. Sets, quantifiers, general recursion, arbitrary code,
currency conversion and unbounded temporal formulas are outside RuleIR 0.1.
An external, evidenced normalizer can supply a reviewed count or HKD value.

These restrictions keep common questions tractable and make the trusted code
small enough to inspect. They also provide a natural place to evaluate the
literature: Catala for the decision fragment, DMN for bank-facing integration,
and a Stipula-inspired state model for a bounded workflow. The project should not
integrate all candidate engines before demonstrating one complete use case.

A reference interpreter implements the written execution semantics. A schema
validator checks structure; a semantic validator checks expression types,
resolved dependencies,
source links and explicit exception relationships; a scenario runner evaluates
reviewed examples. A solver adapter then searches for conflicts, unreachable
branches and semantic differences. Generated code should come from a deterministic
compiler for the supported fragment. At transaction time, evaluation uses an
approved rule version and a recorded fact snapshot; language-model generation
belongs to authoring and review.

An implementation adapter is accepted only after its treatment of missing
values, numeric boundaries, priority, dates and obligations agrees with the
reference semantics on the required checks. DMN's hit-policy choices and Rego's
undefined results make this comparison substantive
\citep{dmn2024,opadocs}. A backend that cannot express a required distinction
should report that limitation instead of approximating it silently.

\subsubsection{Source identity, review and storage}

The source store preserves original bytes and an exact text derivative. The raw
document and derivative each have a SHA-256 digest. A source span records the
derivative hash, page, half-open character offsets and a hash of the quoted text.
Extraction changes create a new derivative; they do not silently move an approved
span. A reviewer can open the original page beside the extracted passage.
Scanned or poorly extracted clauses require visual confirmation before release.

The source register distinguishes circulars, annexes, FAQs, codes, bank policy
and interpretation. Dependency rows record references, definitions, amendments
and replacements; unresolved targets remain explicit. Publication date does not
automatically establish effectiveness. Applicability records identify legal
entity, regulated role, activity, product, client class and jurisdiction. A
product-provider requirement cannot enter the bank's distributor rules merely
because its title mentions a product the bank sells.

The bundle contains definitions, expressions, spans and interpretations. Its
canonical JSON hash identifies the entire immutable revision. Review decisions
live outside those bytes and approve that exact hash. The lifecycle is draft,
in review, approved, released and retired; rejection returns a revision to draft,
while an edit creates a new draft hash. Release requires distinct compliance and
engineering reviewers, resolved blocking questions and dependencies, a supported
effective interval and passing evidence for that hash. The author cannot approve
their own revision under the proposed prototype policy.

SQLite transactions enforce review revisions and idempotency. A stale review
update is rejected rather than overwriting another review. Content-addressed
blobs are written before their database references; a missing referenced blob
is an integrity error. Audit rows record accepted transitions. Hashes detect
accidental corruption but do not make a database tamper-proof against an
administrator able to rewrite both records and hashes. Enterprise identity,
signed retention and access controls belong to the later bank deployment.

\subsubsection{A source-to-decision walkthrough}

The first demonstration imports 23EC35 and Annex 1, verifies their hashes and
opens paragraph 3.1 beside the financial rule. Its two wealth definitions are
review issues until the relevant source and data mappings are accepted. An
engineer loads the supplied draft bundle and the synthetic \texttt{f01} case:
portfolio HK\$40 million and qualifying net assets zero. The trace evaluates
both inclusive comparisons and their alternative, returning \texttt{TRUE} for
\texttt{spi.financial}. It also identifies the exact source revision, input
snapshot and engine version.

The reviewer next opens \texttt{f02}, one cent below the portfolio threshold,
with the other route still false. This returns \texttt{FALSE}. In \texttt{f07},
the portfolio is unknown and net assets are HK\$70 million, so the answer is
\texttt{UNKNOWN}. In \texttt{f06}, changing only net assets to HK\$90 million
establishes the financial condition despite the still-missing portfolio. These
cases let a compliance reader inspect the interpretation before any system rule
is exported. A proposed strict comparison is challenged with \texttt{f01}.

Once dependencies and meanings are reviewed, the evidence report is attached to
the bundle hash, reviewers approve it and the release service exports the package.
The initial production adapter is a deterministic decision API; additional bank
configuration or code generation is a later adapter against the same contract.
A changed rule creates another hash and invalidates release approval. Historic
decisions still point to the old hash and their original evidence.

\subsubsection{Service boundaries an implementer can use}

The API imports sources, creates and retrieves bundles, records reviews, releases
approved hashes, evaluates snapshots, appends events, replays histories and
compares versions. The full method/path/payload table is supplied in the contract
file. The central call is
\begin{lstlisting}
evaluate(bundle, snapshot, rule_id, valid_at, known_at,
         mode="draft")
    -> EvaluationResult
\end{lstlisting}
Its response includes the tagged outcome, exact value if any, reason codes,
missing and blocking inputs, trace, source/bundle/input hashes and engine version.
Production calls require a released exact hash; draft calls visibly retain draft
mode. An HTTP failure is different from a well-formed request whose assessment
returns an error or conflict.

Longer solver or drafting calls use persisted jobs with explicit failed,
cancelled and timed-out states. They cannot alter approval records. Official
documents and model output remain untrusted input: fetching restricts hosts and
redirects, validates content and limits resources; parsing never executes a
document's instructions. The initial service binds to localhost and uses public
sources and synthetic facts. This makes the proposed behavior implementable
before private data access or production integration is considered.


\subsection{Language-model drafting with visible uncertainty}

The drafting service should propose clause segments, roles, definitions,
cross-references, candidate rules and review questions. Every proposal must cite
the exact source passage and expose the definitions it imports. The service
should use constrained structured output, followed by validation against the
rule grammar. Retrieved text is evidence to interpret, not an instruction that
can modify the workbench or execute arbitrary tools.

Independent candidate translations can help discover uncertainty. Their value
comes from meaningful differences, not a vote count. When two candidates disagree,
the solver should try to construct a case that distinguishes them. When both
depend on a missing FAQ, the next action is to retrieve and inspect that source.
When the language requires legal judgment, the result is an explicit assessment
task. The model's confidence and the solver's ability to execute a formula must
remain separate from approval of the interpretation, consistent with the lessons
of ARc and the tax-reasoning work \citep{arc2025,jurayj2026}.

The reviewer should receive a short set of consequential questions. For example:
does the selected evidence establish that an agreement is absent, or merely that
it has not been retrieved? Does a new circular replace the entire earlier source
or a specified part? Does a control apply to the distributor or the provider?
Question selection should reduce ambiguity in the actual decision rather than
produce generic compliance checklists.

\subsection{Amendments and the release package}

Every source revision should produce both a textual change view and an impact
view. The former identifies changed passages, annexes and links; the latter
follows dependencies to interpretations, definitions, rules, tests, data mappings
and deployed controls. The 2023 and 2026 tokenisation circulars form a useful
initial replacement case. A replacement marker is evidence about the document
relationship, but application dates and treatment of existing arrangements still
require the relevant source and review \citep{sfctoken2023,sfctoken2026}.

Where old and new executable rules cover the same domain, the workbench should
find representative cases on which their outcomes differ. Where a new rule
introduces an additional fact or actor, it should identify the new data dependency
and the cases that cannot yet be evaluated. A textual difference alone cannot
measure operational impact, and an unchanged threshold cannot establish that a
definition or exception is unchanged.

The release package should contain the source snapshot, interpretation record,
rule version, unresolved exclusions, approved examples, test and solver results,
adapter version, data mappings and owner approvals. Decisions made with that
package should record its identifier and input snapshot. Dependency analysis can
identify tests and interpretation evidence that remain relevant after a change.
Every changed bundle nevertheless needs fresh approval of its exact hash and a
new evidence report; a result for an older hash cannot authorize its release.
The historical record remains intact. This distinction permits reuse of work
without obscuring what needs fresh review.
